\documentclass[crop=false]{standalone}
\usepackage{15150toolbox}
\usepackage{comment}
\usepackage{needspace}

% Toggle: set \soltrue to render solutions inline (blue boxes).
\newif\ifsol
\soltrue   % comment this line for the student version

% Solution environment.
\ifsol
  \newenvironment{solution}
    {\par\color{solution_color}\begin{framed}\textbf{Solution.}\par}
    {\end{framed}}
\else
  \excludecomment{solution}
\fi

% Lecture-slide macros not in 15150toolbox.
\newcommand{\term}[1]{\textbf{\textit{#1}}}
\newcommand{\thmBox}[2]{\par\noindent\begin{mdframed}[backgroundcolor=blue!8,hidealllines=true]\textbf{Theorem#1.} #2\end{mdframed}}
\newcommand{\bcBox}[1]{\par\noindent\textbf{Base case.} #1\par}
\newcommand{\ihBox}[2]{\par\noindent\textbf{Inductive hypothesis#1.} #2\par}
\newcommand{\isBox}[2]{\par\noindent\textbf{Inductive step#1.} #2\par}

% Problem header.
\newcommand{\quizproblem}[3]{%
  \needspace{5\baselineskip}%
  \stepcounter{problem}%
  \ifsol
    \subsection*{Problem \arabic{problem}. #1 \hfill \timing{#2} \quad (#3 pts)}%
  \else
    \subsection*{Problem \arabic{problem}. #1 \hfill (#3 pts)}%
  \fi
}

\newcommand{\subpts}[1]{\textbf{(#1 pts)}}

\newcommand{\answerspace}[1]{\ifsol\else\vspace{#1}\par\fi}

\printout{Quiz 4}

\begin{document}

\begin{center}
  {\Large \textbf{15-150 --- Quiz 4}} \\
  \vspace{4pt}
  Topics: \textit{Continuation-Passing Style}; \textit{Exceptions};
  \textit{Regular Expressions}; \textit{Modules I} \\
  \vspace{4pt}
  \textit{60 minutes.}
\end{center}

\vspace{1em}

\noindent\textbf{Name:} \hrulefill \hspace{2em} \textbf{Andrew ID:} \hrulefill

\vspace{1em}

% =============================================================================
% PROBLEM 1 --- CPS mechanical translation
% =============================================================================

\quizproblem{Translating to CPS}{14}{14}

Recall that the difference between a continuation-passing-style function and
its direct-style analogue is merely a difference in \emph{contract}. Consider a
function \code{maxOpt : int list -> int option} returning the largest
element of a list, and its CPS version
\code{maxOptCPS : int list -> (int -> 'a) -> (unit -> 'a) -> 'a}. Their
behavior is given by the following specifications:

\[
  \code{maxOpt L} \eeq
  \begin{cases}
    \code{SOME x}, & \text{if the largest element of \code{L} is \code{x}} \\
    \code{NONE},   & \text{if \code{L} is empty}
  \end{cases}
\]
\[
  \code{maxOptCPS L sc fc} \eeq
  \begin{cases}
    \code{sc x},  & \text{if the largest element of \code{L} is \code{x}} \\
    \code{fc ()}, & \text{if \code{L} is empty}
  \end{cases}
\]

\vspace{0.5em}

\textbf{(a)} \subpts{4} The same idea works for arbitrary types.
Consider \code{cmp}, which returns a value of type \code{order}:

\begin{codeblock}
  datatype order = LESS | EQUAL | GREATER

  fun cmp (x : int, y : int) : order =
    if x < y then LESS
    else if x = y then EQUAL
    else GREATER
\end{codeblock}

We want to write \code{cmpCPS}, the CPS version of \code{cmp}. \textbf{Write
a possible type for \code{cmpCPS}, and give its specification in the same
piecewise form as \code{maxOptCPS} above.} (You do not need to implement
it)

Name the three continuations \code{lc}, \code{ec}, and \code{gc}, or
alternatively, \code{snap}, \code{crackle}, and \code{pop}.

\vspace{0.3em}
\hspace{2em}\code{cmpCPS :} \rule[-0.4ex]{0.8\textwidth}{0.4pt}

\answerspace{4cm}

\newpage

\begin{solution}
Type:
\begin{quote}
\code{cmpCPS : int * int ->} \\
\hspace*{2em}\code{(unit -> 'a) -> (unit -> 'a) -> (unit -> 'a) -> 'a}
\end{quote}

Specification:
\[
  \code{cmpCPS (x, y) lc ec gc} \eeq
  \begin{cases}
    \code{lc ()}, & \text{if \code{cmp (x, y)} $\eeq$ \code{LESS}} \\
    \code{ec ()}, & \text{if \code{cmp (x, y)} $\eeq$ \code{EQUAL}} \\
    \code{gc ()}, & \text{if \code{cmp (x, y)} $\eeq$ \code{GREATER}}
  \end{cases}
\]

\textbf{Rubric:}
\begin{itemize}
  \item 2 pts: type has three \code{unit -> 'a} continuation arguments and
    return type \code{'a} (currying the \code{int * int} is fine)
  \item 2 pts: specification maps each \code{order} result to the
    corresponding continuation call (\code{LESS} $\to$ \code{lc ()}, etc.)
\end{itemize}
\end{solution}

\vspace{0.5em}

\textbf{(b)} \subpts{8} Consider the function \code{indexOf}, which returns
the index of the first occurrence of \code{y} in a list, if present:

\begin{codeblock}
  fun indexOf (y : int, [] : int list) : int option = NONE
    | indexOf (y, x::xs) =
        if y = x then SOME 0
        else case indexOf (y, xs) of
               NONE   => NONE
             | SOME i => SOME (i + 1)
\end{codeblock}

Write \code{indexOfCPS}, the CPS version of \code{indexOf}, satisfying

\spec
  {indexOfCPS}
  {int * int list -> (int -> 'a) -> (unit -> 'a) -> 'a}
  {\code{sc} and \code{fc} are total}
  {\vspace{-1.5em}
   {\renewcommand{\arraystretch}{0.9}%
   \[
     \code{indexOfCPS (y, L) sc fc} \eeq
     \begin{cases}
       \code{sc i},  & \text{if \code{indexOf (y, L)} $\eeq$ \code{SOME i}} \\
       \code{fc ()}, & \text{if \code{indexOf (y, L)} $\eeq$ \code{NONE}}
     \end{cases}
   \]}\vspace{-1em}}

\begin{constraint}
  \code{indexOfCPS} must be in proper continuation-passing style.
\end{constraint}

\answerspace{4cm}

\begin{solution}
\begin{codeblock}
  fun indexOfCPS (y, []) sc fc = fc ()
    | indexOfCPS (y, x::xs) sc fc =
        if y = x then sc 0
        else indexOfCPS (y, xs) (fn i => sc (i + 1)) fc
\end{codeblock}

\textbf{Rubric:}
\begin{itemize}
  \item 3 pts: base case calls \code{fc ()}; match case calls \code{sc 0}
  \item 4 pts: recursive case wraps the success continuation ---
    \code{(fn i => sc (i + 1))} --- to perform the \code{+ 1}
  \item 1 pt: recursive call is a tail call, passing \code{fc} unchanged
\end{itemize}
\end{solution}

\vspace{0.5em}

\textbf{(c)} \subpts{2} A student wants to use \code{indexOfCPS} to write
\code{indexOrMinusOne}, which returns the index of \code{y} in \code{L}, or
\code{~1} if \code{y} is absent. They write:

\begin{codeblock}
  fun indexOrMinusOne (y : int, L : int list) : int =
    case indexOfCPS (y, L) (fn i => SOME i) (fn () => NONE) of
      SOME i => i
    | NONE   => ~1
\end{codeblock}

This produces the right answer, but it uses \code{indexOfCPS} incorrectly.
In one sentence, say what's wrong with it, then write the corrected
\code{indexOrMinusOne} below.

\answerspace{4.5cm}

\begin{solution}
\textbf{What's wrong:} This is not CPS!

\textbf{Corrected:}
\begin{codeblock}
  fun indexOrMinusOne (y : int, L : int list) : int =
    indexOfCPS (y, L) (fn i => i) (fn () => ~1)
\end{codeblock}

\textbf{Rubric:}
\begin{itemize}
  \item 1 pt: identifies that the continuations just rebuild an option
    (result-handling should be in the continuations, not a \code{case})
  \item 1 pt: correct rewrite --- \code{(fn i => i)} as \code{sc},
    \code{(fn () => ~1)} as \code{fc}
\end{itemize}
\end{solution}

% =============================================================================
% PROBLEM 2 --- CPS implementation from scratch
% =============================================================================

\quizproblem{The Depth of CPS}{10}{12}

Recall the type of binary trees:

\begin{codeblock}
  datatype tree = Empty | Node of tree * int * tree
\end{codeblock}

The \code{depth} of a tree is the length of the longest path from the root
to a leaf. In direct style:

\begin{codeblock}
  fun depth (Empty : tree) : int = 0
    | depth (Node (L, _, R)) = 1 + Int.max (depth L, depth R)
\end{codeblock}

Write \code{depthCPS}, a continuation-passing-style version of \code{depth},
satisfying

\spec
  {depthCPS}
  {tree -> (int -> 'a) -> 'a}
  {\code{k} is total}
  {\code{depthCPS T k} $\eeq$ \code{k (depth T)}}

\begin{constraint}
  \code{depthCPS} must be in proper continuation-passing style.
\end{constraint}

You may assume \code{Int.max} is given.

\answerspace{6cm}

\begin{solution}
\begin{codeblock}
  fun depthCPS (Empty : tree) (k : int -> 'a) : 'a = k 0
    | depthCPS (Node (L, _, R)) k =
        depthCPS L (fn dl =>
          depthCPS R (fn dr =>
            k (1 + Int.max (dl, dr))))
\end{codeblock}

\textbf{Rubric:}
\begin{itemize}
  \item 2 pts: base case \code{k 0}
  \item 5 pts: both subtrees recursed on, with the second call nested
    inside the first's continuation (not two separate top-level calls)
  \item 3 pts: \code{k} called only at the bottom, on
    \code{1 + Int.max (dl, dr)}
  \item 2 pts: every call is a tail call (proper CPS, no pending work after
    a recursive call)
\end{itemize}
\end{solution}

\newpage

% =============================================================================
% PROBLEM 3 --- CPS with a two-continuation predicate
% =============================================================================

\quizproblem{If I Had To Count The Number Of Things That Were Cool About
Functional Programming, I Would Have To Count To Quite A High Number Indeed}{9}{8}

We are interested in predicates that can be in CPS form too. Instead of
\code{p : int -> bool}, consider a predicate

\begin{codeblock}
  p : int -> (unit -> 'a) -> (unit -> 'a) -> 'a
\end{codeblock}

where \code{p x yes no} calls \code{yes ()} if \code{x} satisfies the
predicate, and \code{no ()} if it does not. (This is as you saw on
the CPS homework)

Write \code{countCPS}, which counts how many elements of a list satisfy
such a predicate and passes the resulting count to a continuation \code{k}:

\spec
  {countCPS}
  {(int -> (unit -> 'a) -> (unit -> 'a) -> 'a) * int list
   -> (int -> 'a) -> 'a}
  {\code{p} and \code{k} are total, \code{p} is a CPS predicate function
  as described above}
  {\code{countCPS (p, L) k} $\eeq$ \code{k n}, where \code{n} is the number
  of elements of \code{L} satisfying \code{p}}

\begin{constraint}
  \code{countCPS} must be in proper continuation-passing style.
\end{constraint}

Complete the two clauses below.

\begin{codeblock}
  fun countCPS (p, []) k =

    | countCPS (p, x::xs) k =
\end{codeblock}

\answerspace{5cm}

\begin{solution}
\begin{codeblock}
  fun countCPS (p, []) k = k 0
    | countCPS (p, x::xs) k =
        p x
          (fn () => countCPS (p, xs) (fn n => k (n + 1)))
          (fn () => countCPS (p, xs) k)
\end{codeblock}

\textbf{Rubric:}
\begin{itemize}
  \item 4 pts: calls \code{p x} with two continuations --- yes (count
    \code{x}), no (skip \code{x})
  \item 2 pts: the ``yes'' branch wraps the continuation as
    \code{(fn n => k (n + 1))} to add \code{1} to the rest's count
  \item 1 pt: the ``no'' branch recurses with \code{k} unchanged
  \item 1 pt: base case calls \code{k 0}; every call is a tail call
\end{itemize}
\end{solution}

\newpage

% =============================================================================
% PROBLEM 4 --- Exceptions
% =============================================================================

\quizproblem{Indexing, Exceptionally}{9}{11}

In \term{exception-handling style} (EHS), a function uses exceptions ---
both raising \emph{and} handling them --- to carry information and direct
control flow. We will write \code{indexOf} (from Problem 1) in EHS: it
raises \code{Found i} if \code{y} occurs at index \code{i}, and
\code{NotFound} if \code{y} is absent. It never returns normally.

\begin{codeblock}
  exception NotFound
  exception Found of int
\end{codeblock}

\textbf{(a)} \subpts{8} Write \code{indexOfEHS}, satisfying

\spec
  {indexOfEHS}
  {int * int list -> 'a}
  {\code{true}}
  {\code{indexOfEHS (y, L)} raises \code{Found i} if the first occurrence of
   \code{y} in \code{L} is at index \code{i}, and raises \code{NotFound} if
   \code{y} does not occur in \code{L}}

\begin{constraint}
  \code{indexOfEHS} must be in exception-handling style (use \code{handle},
  not a helper that returns an \code{int} or \code{option}).
\end{constraint}

\answerspace{6cm}

\begin{solution}
\begin{codeblock}
  fun indexOfEHS (y : int, [] : int list) : 'a = raise NotFound
    | indexOfEHS (y, x::xs) =
        if y = x then raise (Found 0)
        else indexOfEHS (y, xs) handle Found i => raise (Found (i + 1))
\end{codeblock}

\textbf{Rubric:}
\begin{itemize}
  \item 3 pts: base case raises \code{NotFound}; match case raises
    \code{Found 0}
  \item 5 pts: recursive case is \code{indexOfEHS (y, xs) handle Found i =>
    raise (Found (i + 1))} --- catches the inner result and re-raises it
    incremented, while letting \code{NotFound} propagate
\end{itemize}
\end{solution}

\vspace{0.5em}

\textbf{(b)} \subpts{3} Since \code{indexOfEHS} never returns, a caller must
\code{handle} its exceptions. Complete \code{indexOrMinusOne}, which
evaluates to the index of \code{y} in \code{L}, or \code{~1} if \code{y} is
absent.

\begin{codeblock}
  fun indexOrMinusOne (y : int, L : int list) : int =
\end{codeblock}

\answerspace{3.5cm}

\begin{solution}
\begin{codeblock}
  fun indexOrMinusOne (y : int, L : int list) : int =
    indexOfEHS (y, L)
      handle Found i => i
           | NotFound => ~1
\end{codeblock}

\textbf{Rubric:}
\begin{itemize}
  \item 2 pts: \code{handle Found i => i} (extracts the index)
  \item 1 pt: \code{handle NotFound => ~1} for the absent case
\end{itemize}
\end{solution}

\newpage

% =============================================================================
% PROBLEM 4 --- Regular Expressions
% =============================================================================

\quizproblem{Match Point}{5}{5}

Consider the alphabet $\Sigma = \{\code{a}, \code{b}\}$ and the regular
expression
$$ r = \code{a}(\code{a} + \code{b})^* \code{b} $$
that is, the language of strings that start with \code{a}, end with
\code{b}, and have any string of \code{a}'s and \code{b}'s in between.

For each string below, circle \textbf{YES} if it is in $L(r)$ and
\textbf{NO} if it is not.

\vspace{0.5em}

\begin{center}
\begin{tabular}{l@{\hspace{3em}}c@{\hspace{2em}}c}
  \code{ba}     & YES \quad NO & \\
  \code{aab}    & YES \quad NO & \\
  \code{a}      & YES \quad NO & \\
  \code{abb}    & YES \quad NO & \\
  \code{ab}     & YES \quad NO & \\
\end{tabular}
\end{center}

\answerspace{1cm}

\begin{solution}
\begin{center}
\begin{tabular}{l@{\hspace{2em}}l@{\hspace{2em}}l}
  \code{ba}  & \textbf{NO}  & does not start with \code{a}. \\
  \code{aab} & \textbf{YES} & \code{a}, then \code{a} from the star, then \code{b}. \\
  \code{a}   & \textbf{NO}  & has no final \code{b} (the leading \code{a} is consumed). \\
  \code{abb} & \textbf{YES} & \code{a}, then \code{b} from the star, then \code{b}. \\
  \code{ab}  & \textbf{YES} & \code{a}, then $\epsilon$ from the star, then \code{b}. \\
\end{tabular}
\end{center}

\textbf{Rubric:} 1 pt each, all-or-nothing.
\end{solution}

% =============================================================================
% PROBLEM 5 --- Modules
% =============================================================================

\quizproblem{Abstract Thinking}{8}{9}

Consider the following signature and two structures that implement it ---
one ascribed \textbf{transparently} (\code{:}) and one ascribed
\textbf{opaquely} (\code{:>}). Both use \code{type t = int} internally.

\begin{codeblock}
  signature COUNTER =
    sig
      type t
      val start : t
      val tick  : t -> t
    end

  structure C : COUNTER =
    struct type t = int  val start = 0  fun tick n = n + 1 end

  structure D :> COUNTER =
    struct type t = int  val start = 0  fun tick n = n + 1 end
\end{codeblock}

For each client expression below, state whether it \textbf{typechecks}, and
explain in one sentence why or why not.

\vspace{0.5em}

\textbf{(a)} \subpts{3} \code{C.tick C.start + 1}

\answerspace{2.5cm}

\textbf{(b)} \subpts{3} \code{D.tick D.start + 1}

\answerspace{2.5cm}

\begin{solution}
\textbf{(a) Typechecks.} \code{C} is ascribed transparently, so
\code{C.t} is publicly known to be \code{int}.

\textbf{(b) Does not typecheck.} \code{D} is ascribed opaquely,
so \code{D.t} is abstract and \code{+}
cannot be applied to a value of type \code{D.t}.

\textbf{Rubric:}
\begin{itemize}
  \item 3 pts (a): typechecks, because transparent ascription exposes
    \code{C.t = int}. (1 pt for ``yes'' with weak/no reason.)
  \item 3 pts (b): does not typecheck, because opaque ascription makes
    \code{D.t} abstract, so it isn't known to be \code{int} and \code{+}
    can't apply. (1 pt for ``no'' with weak/no reason.)
\end{itemize}
\end{solution}

\vspace{0.5em}

\textbf{(c)} \subpts{3} Now consider a third structure that \emph{reuses
\code{C}'s type}, ascribed opaquely:

\begin{codeblock}
  structure E :> COUNTER =
    struct type t = C.t  val start = C.start  fun tick n = C.tick n end
\end{codeblock}

Does the client expression \code{C.start = E.start} typecheck? Why or why
not?

\answerspace{3cm}

\begin{solution}
\textbf{Does not typecheck.} Even though \code{E} reuses \code{C.t}
internally, the opaque ascription does not let us know about it. So
the equality fails to type-check.

\textbf{Rubric:}
\begin{itemize}
  \item 3 pts: does not typecheck, because \code{E.t} is a fresh abstract
    type distinct from \code{C.t} (so \code{=} compares two different
    types). (1 pt for ``no'' with weak/no reason, e.g. only mentioning
    that the type is hidden or equality types.)
\end{itemize}
\end{solution}

\end{document}
